\section{Introduction}
La distance globale compte les éditions nécessaires pour transformer une chaîne entière en une autre. NvPD, proposé par Sadiq et ses coauteurs, calcule les colonnes de chaque ligne en parallèle~\cite{sadiq}. Les essais comparent son travail dense sur EPYC et GH200 aux méthodes qui dépendent du score.

Pour deux longueurs \(m,n\) et des coûts unitaires,
\[
 |m-n|\le d(A,B)\le\max(m,n),\qquad W_{\mathrm{dense}}=mn.
\]
\synthesesection{Deux chaînes d'un million de bases définissent près de \(10^{12}\) cases. Les méthodes denses les visitent toutes, tandis que le nombre de fronts nécessaires dépend du score atteint.}

\section{Temps d'une paire et débit d'une collection}
\subsection{Longueur et divergence}
Edlib de Šošić et Šikić code plusieurs différences de scores dans un mot binaire~\cite{edlib}. WFA2 de Marco-Sola et ses coauteurs garde la position la plus avancée sur chaque diagonale à coût fixé~\cite{wfa}. Sur GH200, Kokkos calcule deux lignes de NvPD par lancement. Pour une fenêtre humaine de 524\,288 bases, cette fusion ramène le temps de \(2{,}520\pm0{,}034\) à \(1{,}345\pm0{,}018\)\,s (cinq processus après échauffement).
\chart{comparison/synthese.png}{Le temps d'une paire est comparé au débit d'un lot de taille fixe. Les barres d'incertitude de chaque panneau portent sur ses unités expérimentales.}{fig:short-work}{notebooks/profil_nvpd.ipynb}
La figure~\ref{fig:short-work} compare les régimes de calcul sur les mêmes jeux de paires. Sur une paire à 2\,\%, cinq blocs chauds donnent \(3{,}570\pm0{,}009\)\,s avec NvPD sur GH200 et \(0{,}767\pm0{,}006\)\,s avec WFA2 sur Grace. À 10\,\%, les mêmes méthodes prennent \(3{,}591\pm0{,}011\) et \(9{,}835\pm0{,}023\)\,s.

\synthese{Les trois taux injectés, 2, 5 et 10\,\%, laissent l'aire du rectangle presque constante. Ils modifient pourtant assez le temps de WFA2 pour inverser son classement face au calcul dense sur GH200.}

\subsection{Communication entre nœuds}
Avec MPI, les rangs peuvent partager une paire en bandes verticales ou traiter des paires indépendantes. Sur une paire humaine de 131\,072 bases et deux nœuds, échanger le bord par paquets de 128 lignes ramène le temps de \(6{,}672\pm0{,}182\)\,s à \(0{,}117\pm0{,}001\)\,s. Sur un lot, les rangs conservent leurs scores jusqu'à la réduction finale. L'efficacité pour un travail fixe vaut \(E=T_{192}/(2T_{384})\). Les temps proviennent de cinq blocs avec un échauffement par configuration.
\synthese{Pour 768 paires homologues du chromosome 21, NvPD atteint 80,8\,\% d'efficacité forte sur 384 cœurs par rapport à 192 cœurs. L'efficacité faible, corrigée du travail réel, vaut 80,9\,\%.}
\synthesesection{À deux nœuds, la décomposition d'une paire nécessite encore des échanges intermédiaires. Un lot de 768 paires permet d'attribuer un travail complet à chacun des 384 cœurs.}

\section{Assemblages humains}
Deux versions de CHM13 diffèrent de 2\,587 éditions sur dix-huit chromosomes entiers, soit 2,644 milliards de bases de référence. Sur les deux entrées du chromosome~1, le score vaut 176. La lecture seule prend \(0{,}657\pm0{,}005\)\,s sur un cœur EPYC avec préallocation par l'index FASTA (neuf processus).
\chart{comparison/bras_chromosome_biwfa.png}{Les temps de calcul et les pics résidents mesurés sur le bras long du chromosome humain dépendent du score et de la restitution du chemin.}{fig:short-chrom}{bench/human/algorithms/arm.py}
La figure~\ref{fig:short-chrom} porte sur GRCh38 et CHM13, deux références distinctes. Les fenêtres de 8 à 32 millions de bases partagent la même borne gauche. BiWFA de Marco-Sola et ses coauteurs propage deux fronts opposés et reconstruit le chemin de part et d'autre de leur jonction optimale~\cite{biwfa}. Sa mémoire de travail est en \(O(d)\) en plus des entrées et du chemin rendu.
\synthesesection{Sur l'intervalle de 32 millions de bases, les 266\,220 éditions demandent \(45{,}62\pm0{,}22\)\,s pour le score BiWFA et \(89{,}77\pm0{,}11\)\,s pour le chemin sur un cœur EPYC.}

\section{Conclusion}
La distance guide le choix des fronts ou du calcul dense. Pour les 64 phages, une fois les scores calculés, l'erreur entre la distance d'édition \(d_{ij}\) et la longueur du chemin dans l'arbre \(t_{ij}\) mesure la perte associée à la reconstruction,
\[
 \sqrt{\frac{1}{2\,016}\sum_{i<j}(d_{ij}-t_{ij})^2}
 = \begin{cases}14{,}89\ \text{éditions} & \text{(UPGMA)},\\
                  3{,}56\ \text{éditions} & \text{(Neighbor-Joining)}.\end{cases}
\]
\synthesesection{Sur les 64 phages, NJ réduit l'erreur quadratique moyenne des chemins, mais son écart absolu maximal atteint 63,86 éditions. Une distance d'édition exacte peut donc être imparfaitement représentée par un arbre.}

\FloatBarrier
\begin{thebibliography}{4}
\bibitem{sadiq} M. U. Sadiq et coll. \href{https://doi.org/10.1007/s10586-019-02962-w}{\emph{NvPD: novel parallel edit distance algorithm, correctness, and performance evaluation.}} Cluster Computing 23 (2020).
\bibitem{edlib} M. Šošić et M. Šikić. \href{https://doi.org/10.1093/bioinformatics/btw753}{\emph{Edlib: a C/C++ library for fast, exact sequence alignment using edit distance.}} Bioinformatics 33 (2017).
\bibitem{wfa} S. Marco-Sola et coll. \href{https://doi.org/10.1093/bioinformatics/btaa777}{\emph{Fast gap-affine pairwise alignment using the wavefront algorithm.}} Bioinformatics 37 (2021).
\bibitem{biwfa} S. Marco-Sola et coll. \href{https://doi.org/10.1093/bioinformatics/btad074}{\emph{Optimal gap-affine alignment in \(O(s)\) space.}} Bioinformatics 39 (2023).
\end{thebibliography}
